\section*{Chapter \ref{ch:g9:first-cryptography} --- Secret Codes: A First Look at Cryptography}

\begin{solution}{exo:g9:first-cryptography:1}
\code{frgh}: c, o, d, e move three places to f, r, g, h.
\end{solution}

\begin{solution}{exo:g9:first-cryptography:2}
Moving back 23 places is the same as moving forward 3: \code{qeb} becomes
\code{the}.
\end{solution}

\begin{solution}{exo:g9:first-cryptography:3}
It multiplies the numbers of choices: $4 \times 3 \times 2 \times 1 = 24$. Four
letters can be mixed in 24 ways, so there are 24 keys, one of which (the
alphabet unchanged) hides nothing.
\end{solution}

\begin{solution}{exo:g9:first-cryptography:4}
Capitals have their own codes, from \code{ord("A")}: the \omterm{def:g7:functions:function}{function} treats each
case with its own starting letter.

\omcode{code/grade-9/06-first-cryptography/exercises/shift_case.py}{1}{17}{A shift
that keeps capitals.}\label{lst:g9:first-cryptography:sol-case}

It prints \code{Phhw Ph dw Qrrq, eb wkh Crr!}.
\end{solution}

\begin{solution}{exo:g9:first-cryptography:5}
It prints \texttt{d \{ \}}: the codes after z's are those of other
\omterm{def:g7:strings:string}{characters}, the braces \texttt{\{} and \texttt{\}}. The fix counts the position in the alphabet and keeps
the remainder by 26, as \code{shift_letter} does in
\cref{lst:g9:first-cryptography:shift}.
\end{solution}

\begin{solution}{exo:g9:first-cryptography:6}
It prints \code{coo iye kd dro rkblyeb}: the message shifted by 10, since it
shifts by 5 a second time instead of back. Decrypt with \code{encrypt(secret, -5)}
(or with \code{21}, which goes round to the same place).
\end{solution}

\begin{solution}{exo:g9:first-cryptography:7}
A \omterm{def:g7:functions:function}{function} scores each candidate by its common words, and the \omterm{def:g4:algorithms-repeat-and-decide:loop}{loop} keeps the
best:

\omcode{code/grade-9/06-first-cryptography/exercises/best_guess.py}{1}{24}{The
\omterm{def:g9:first-cryptography:encryption}{decryption} with the most common words.}\label{lst:g9:first-cryptography:sol-best}

It prints \code{the keeper lit the lamp at the end of the day} (the key was 2).
\end{solution}

\begin{solution}{exo:g9:first-cryptography:8}
The same \omterm{def:g4:algorithms-repeat-and-decide:loop}{loop} as \code{keyword_encrypt}, with each shift taken the other way:

\omcode{code/grade-9/06-first-cryptography/exercises/keyword_decrypt.py}{1}{18}{Undoing
the \omterm{def:g9:first-cryptography:keyword-cipher}{keyword cipher}.}\label{lst:g9:first-cryptography:sol-keyword}

It prints \code{meet me at the lighthouse at nine}.
\end{solution}

\begin{solution}{exo:g9:first-cryptography:9}
A shift of 26 takes every letter all the way round, back to itself: the message
does not change. A shift of 29 is the same as a shift of 3, because only the
remainder by 26 matters. So the keys 1 to 25 are the only different ones that
change the message.
\end{solution}

\begin{solution}{exo:g9:first-cryptography:10}
\omcode{code/grade-9/06-first-cryptography/exercises/all_keys_time.py}{1}{6}{Trying
every substitution key.}\label{lst:g9:first-cryptography:sol-time}

It prints \code{12,788,288,341 years}: almost thirteen billion years. Trying
every key is hopeless; \omterm{def:g9:first-cryptography:frequency-analysis}{frequency analysis} does not need to.
\end{solution}

\begin{solution}{exo:g9:first-cryptography:11}
\omcode{code/grade-9/06-first-cryptography/exercises/short_counts.py}{1}{11}{Counts on
the start of the passage.}\label{lst:g9:first-cryptography:sol-short}

It prints, for 20 letters, \code{t 3}, \code{h 3}, \code{e 2}, \code{o 2},
\code{l 2}; for 100, \code{e 18}, \code{t 11}, \code{h 11}, \code{r 9},
\code{a 8}; for 1,000, \code{e 150}, \code{t 123}, \code{h 80}, \code{o 68},
\code{a 68}. In 20 letters, every letter appears only a few times, and one word
more or less changes the order: e is not even first. With 1,000 letters, the
order is close to that of the whole passage. A frequency guess needs enough
letters for the counts to settle, as the frequency of a six did in
\cref{ch:g9:randomness-and-simulation}.
\end{solution}

\begin{solution}{exo:g9:first-cryptography:12}
A \omterm{def:g8:dictionaries:dictionary}{dictionary} remembers where each group of three letters was last seen:

\omcode{code/grade-9/06-first-cryptography/exercises/repeats.py}{1}{23}{Repeated
groups and their distances.}\label{lst:g9:first-cryptography:sol-repeats}

It finds 372 repeats. Among the lengths 3 to 8, length 4 divides by far the most
distances: 333, against 132 for 3 and 155 for 8. A word such as \code{the} gives
the same three \omterm{def:g9:first-cryptography:cipher}{ciphertext} letters only when it falls under the same part of the
keyword, so the distance between the two places is a multiple of 4. Length 2
divides even more distances (352), but only because every multiple of 4 is also
even: the keyword's length is the largest length that divides nearly as many.
\end{solution}

\begin{solution}{pb:g9:first-cryptography:1}
\textbf{1.} \code{z} appears 9 times, \code{p} 8, \code{l} 7, \code{e} 6.

\textbf{2.} Counting from 0, as \code{ord(letter) - ord("a")} does, z is at
position 25 and e at position 4: $25 - 4 = 21$.

\textbf{3.} With 21, \code{esp} becomes \code{jxu}, which means nothing. With the
second letter, p: $15 - 4 = 11$, and \code{esp} becomes \code{the}. On one line,
o (here z) happened to beat e.

\textbf{4.} That the most frequent letter of the \omterm{def:g9:first-cryptography:cipher}{ciphertext} stands for e. The key
is the distance from e to that letter.

\textbf{5.} The letter may come before e in the alphabet: for b, $1 - 4 = -3$, and
the remainder by 26 turns it into the key 23.

\omoutput{code/grade-9/06-first-cryptography/outputs/guess_shift.txt}{The whole
paragraph, decrypted.}\label{out:g9:first-cryptography:guess-shift}

\textbf{6.} Yes: the key 11, and every letter of the paragraph is right.

\textbf{7.} On a message where e is not the most frequent letter: a \omterm{def:g7:lists:list}{list} of
names, a text in another language, or a text written to avoid the letter e.

\textbf{8.} Each message is a run of letters cut from a random place of the
passage, and each is encrypted with its own random key from 1 to 25.

\omoutput{code/grade-9/06-first-cryptography/outputs/short_messages.txt}{The share
of keys found, by length.}\label{out:g9:first-cryptography:short}

\textbf{9.} 23 times out of 100 with 10 letters, 56 with 100.

\textbf{10.} In this passage, t comes close behind e (11.9\,\% against 14.0\,\%):
in 100 or 200 letters, a few more t's often put t first. With 500 letters, the
gap of about ten letters is larger than the luck of the cut. And runs of 500 or
1,000 letters take a large part of a 1,517-letter passage, where e leads
clearly.

\textbf{11.} A short note is hard to break by counting: its frequencies are not
yet those of the language. But the \omterm{def:g9:first-cryptography:shift-cipher}{shift cipher} has only 25 keys, so trying
them all still breaks any note, however short.

\textbf{12.} \textbf{Result.} The guess from the most frequent letter finds the key
of a 10-letter message 23 times out of 100, of a 100-letter one 56 times, and of
messages of 500 or 1,000 letters every time.
\end{solution}
